\SGANextSectionTarget{sga3:xxv:subsection:2-6}\SGASubsection{2.6}{}
\SGARefTarget{sga3:xxv:paragraph:2:6}

Replacing \(R^+\) by \(R^-=-R^+\), one similarly constructs the groups
\(U^-\), \(B^-\), and \(U_\alpha\) for \(\alpha\in R^-\), together with
isomorphisms
\[
 p_\alpha\colon\mathbb G_a\xrightarrow{\ \sim\ }U_\alpha,
 \qquad \alpha\in R^-.
\]
Finally introduce the product scheme
\[
 \Omega=U^-\times T\times U.
\]
There is a canonical isomorphism of \(k\)-schemes
\[
 \Omega\otimes_{\mathbb Z}k
 \xrightarrow{\ \sim\ }
 U^-_k\times_kT_k\times_kU_k
 \simeq\Omega_k,
\]
where \(\Omega_k\) is the ``big cell'' of \(G_k\) (\SGARefLink{sga3:xxii:prop:4.1.2}{Exposé XXII, 4.1.2}).

From now on we identify \(\Omega\otimes_{\mathbb Z}k\) with \(\Omega_k\)
by this isomorphism, and regard \(U^-\), \(T\), and \(U\) as subschemes
of \(\Omega\) through their unit sections. Write
\[
 e=((0),e,(0))
\]
for the ``unit section'' of \(\Omega\). Our aim is now to put a
piece-of-group law on \(\Omega\).

\begin{sgastatement}{\SGAInlineRefTarget{sga3:xxv:lemma:2-7}Lemma 2.7}
Let \(\alpha\in\Delta\), and let \(w_{\alpha,k}\) be the element of
\(\operatorname{Norm}_{G_k}(T_k)(k)\) defined by \(X_{\alpha,k}\); recall
that, by definition,
\[
 w_{\alpha,k}
 =p_{-\alpha,k}(-1)p_{\alpha,k}(1)p_{-\alpha,k}(-1).
\]
There exist an open subset \(V_\alpha\) of \(\Omega\), containing the
section \(e\), and a morphism
\[
 h_\alpha\colon V_\alpha\longrightarrow\Omega
\]
such that:
\begin{enumerate}[label=\textup{(\roman*)},leftmargin=2.2em]
\item \(h_\alpha(e)=e\);
\item \((h_\alpha)\otimes_{\mathbb Z}k\) coincides with the restriction
of \(\operatorname{int}(w_{\alpha,k})\) to
\(V_\alpha\otimes_{\mathbb Z}k\subset G_k\);
\item \(T\subset V_\alpha\), and \(h_\alpha\) sends \(T\) into \(T\).
For every \(\beta\in R\), one has \(U_\beta\subset V_\alpha\), and
\(h_\alpha\) sends \(U_\beta\) into \(U_{s_\alpha(\beta)}\).
\end{enumerate}
\end{sgastatement}
\SGARefTarget{sga3:xxv:lemma:2:7}

By the definition of a Chevalley system (\SGARefLink{sga3:XXIII:6.1}{Exposé XXIII, 6.1}), for every
\(\beta\in R\) there exists an integer \(e_\beta=\pm1\) such that
\[
 \operatorname{int}(w_{\alpha,k})p_{\beta,k}(x)
 =p_{s_\alpha(\beta),k}(e_\beta x)
\]
for every \(x\in\mathbb G_a(S)\) and every \(S\to\operatorname{Spec}(k)\).

Let \(S\) be any scheme. Write an arbitrary element of \(\Omega(S)\) in
the form
\[
\left(
 \left(\prod_{\substack{\beta\in R^-\\\beta\ne-\alpha}}
 p_\beta(x_\beta)\right)p_{-\alpha}(x_{-\alpha}),
 \ t,\
 p_\alpha(x_\alpha)
 \left(\prod_{\substack{\beta\in R^+\\\beta\ne\alpha}}
 p_\beta(x_\beta)\right)
\right),
\]
where arbitrary orders have been chosen on \(R^-\setminus\{-\alpha\}\)
and \(R^+\setminus\{\alpha\}\), as permitted by \SGARefLink{sga3:xxv:subsection:2-5}{2.5}. Define a morphism
\(d\colon\Omega\to\mathbb A^1_{\mathbb Z}\) by
\[
 d(u)=\alpha(t)+x_\alpha x_{-\alpha}.
\]

Let \(V_\alpha=\Omega_d\), the open subset of \(\Omega\) defined by the
invertibility of \(d(u)\). It contains \(e\), \(T\), and every \(U_\beta\),
\(\beta\in R\). Define
\[
 h_\alpha\colon V_\alpha\longrightarrow\Omega,\qquad
 h_\alpha(u)=(a(u),b(u),c(u)),
\]
where
\begin{align*}
a(u)
 &=
\left(\prod_{\substack{\beta\in R^-\\\beta\ne-\alpha}}
 p_{s_\alpha(\beta)}(e_\beta x_\beta)\right)
p_{-\alpha}\!\left(-x_\alpha d(u)^{-1}\right),\\
b(u)&=t\cdot\alpha^\ast(d(u)),\\
c(u)
 &=p_\alpha\!\left(-x_{-\alpha}d(u)^{-1}\right)
\left(\prod_{\substack{\beta\in R^+\\\beta\ne\alpha}}
 p_{s_\alpha(\beta)}(e_\beta x_\beta)\right).
\end{align*}
Since \(s_\alpha\) permutes the positive roots distinct from \(\alpha\)
(respectively the negative roots distinct from \(-\alpha\)), \(c(u)\)
(respectively \(a(u)\)) is a section of \(U\) (respectively \(U^-\)),
so the preceding morphism is well defined. It plainly satisfies (i)
and (iii). Assertion (ii) follows immediately from the definition of
the \(e_\beta\), \(\beta\in R\), and \SGARefLink{sga3:xx:3.12}{Exposé XX, 3.12}.

\begin{sgastatement}{\SGAInlineRefTarget{sga3:xxv:lemma:2-8}Lemma 2.8}
There exist open subsets \(V,V'\) of \(\Omega\) and morphisms
\[
 h\colon V\longrightarrow\Omega,\qquad
 h'\colon V'\longrightarrow\Omega
\]
with the following properties:
\begin{enumerate}[label=\textup{(\roman*)},leftmargin=2.2em]
\item \(V\) and \(V'\) contain \(e\), and \(h(e)=h'(e)=e\);
\item\SGAInlineRefTarget{sga3:xxv:lemma:2-8:item:ii} the morphism induced by
\[
 h'\circ h\colon h^{-1}(V')\longrightarrow\Omega
\]
is the restriction of the identity morphism;
\item\SGAInlineRefTarget{sga3:xxv:lemma:2-8:item:iii} the \(k\)-morphisms \(h\otimes_{\mathbb Z}k\) and
\(h'\otimes_{\mathbb Z}k\) are the restrictions, to
\(V\otimes_{\mathbb Z}k\) and \(V'\otimes_{\mathbb Z}k\), of
automorphisms of the group \(G_k\);
\item\SGAInlineRefTarget{sga3:xxv:lemma:2-8:item:iv} \(V\) and \(V'\) contain \(U,T,U^-\); the morphisms \(h,h'\)
send \(U\) into \(U^-\), \(U^-\) into \(U\), and \(T\) into \(T\).
\end{enumerate}
\end{sgastatement}
\SGARefTarget{sga3:xxv:lemma:2:8}

Let \(w^0\) be the element of the Weyl group of \(G_k\) that sends \(R^+\)
to \(R^-\). Write
\[
 w^0=s_{\alpha_n}\cdots s_{\alpha_1},\qquad \alpha_i\in\Delta
\]
(with no relation to the numbering in \SGARefLink{sga3:xxv:subsection:2-1}{2.1}), and put
\[
 w_0=w_{\alpha_n,k}\cdots w_{\alpha_1,k}
 \in\operatorname{Norm}_{G_k}(T_k)(k).
\]
Define recursively, for \(i\leq n\), an open subset \(V_i\) of \(\Omega\)
and a morphism \(h_i\colon V_i\to\Omega\) by
\[
 V_0=\Omega,\qquad h_0=\operatorname{id},
\]
and, for \(i=0,\ldots,n-1\),
\[
 V_{i+1}=h_i^{-1}(V_{\alpha_{i+1}}),\qquad
 h_{i+1}=h_{\alpha_{i+1}}\circ h_i,
\]
where \(V_{\alpha_j}\) and \(h_{\alpha_j}\) are those of \SGARefLink{sga3:xxv:lemma:2-7}{2.7}.

Take \(V=V_n\) and \(h=h_n\). The conditions (i), (iii), and (iv) concerning
\(V\) and \(h\) hold: for (i) and (iii) this follows immediately from
\SGARefLink{sga3:xxv:lemma:2-8}{2.8}, and for (iv) from the fact that \(h\otimes_{\mathbb Z}k\) is the
restriction of \(\operatorname{int}(w_0)\) to
\(V\otimes_{\mathbb Z}k\).

Since \((w^0)^2=1\), we also have
\[
 w^0=s_{\alpha_1}\cdots s_{\alpha_n}
 =(s_{\alpha_1}^{\,3})\cdots(s_{\alpha_n}^{\,3}).
\]
Put
\[
 w'_0=(w_{\alpha_1,k})^3\cdots(w_{\alpha_n,k})^3.
\]
Repeating the construction above gives an open subset \(V'\) and a
morphism \(h'\) that also satisfy (i), (iii), and (iv). Moreover,
\(h'\otimes_{\mathbb Z}k\) is the restriction of
\(\operatorname{int}(w'_0)\) to \(V'\otimes_{\mathbb Z}k\). For every
simple root \(\alpha\in\Delta\), however,
\((w_{\alpha,k})^4=e\) (\SGARefLink{sga3:xx:3.1}{Exposé XX, 3.1}), hence \(w'_0w_0=e\).
It follows that \(h'\circ h\) induces the identity morphism on a nonempty
open subset of \(\Omega\otimes_{\mathbb Z}k\). Since \(\Omega\) is smooth
and of finite presentation over \(\mathbb Z\),
\(\Omega\otimes_{\mathbb Z}k\) is schematically dense in \(\Omega\);
this proves (ii).

\begin{sgastatement}{\SGAInlineRefTarget{sga3:xxv:proposition:2-9}Proposition 2.9}
There exist an open subset \(V_1\) of \(\Omega\times\Omega\), an open
subset \(V_2\) of \(\Omega\), and morphisms
\[
 \pi\colon V_1\longrightarrow\Omega,\qquad
 \sigma\colon V_2\longrightarrow\Omega
\]
having the following properties:
\begin{enumerate}[label=\textup{(\roman*)},leftmargin=2.2em]
\item if \(x\in\Omega(S)\), then \((e,x)\) and \((x,e)\) are sections
of \(V_1\), and
\[
 \pi(e,x)=\pi(x,e)=x;
\]
\item \(V_2\) contains \(e\), and \(\sigma(e)=e\);
\item\SGAInlineRefTarget{sga3:xxv:proposition:2-9:item:iii} \(\pi_k\) and \(\sigma_k\) are the restrictions of the morphisms
\(G_k\times_kG_k\to G_k\) and \(G_k\to G_k\) defined respectively by
multiplication and inversion.
\end{enumerate}
\end{sgastatement}
\SGARefTarget{sga3:xxv:proposition:2:9}

\SGAProof
Let \((v,t,u)\) and \((v',t',u')\) be two sections of \(\Omega\).
By \SGARefLink{sga3:xxv:lemma:2-8:item:iv}{2.8(iv)}, \(h(u)\) is a section of \(U^-\) and \(h(v')\) a section
of \(U\); thus \((h(u),e,h(v'))\) may be regarded as a section of
\(\Omega\). Define \(V_1\) as the open subset of \(\Omega\times\Omega\)
specified by
\[
\bigl((v,t,u),(v',t',u')\bigr)\in V_1(S)
\quad\Longleftrightarrow\quad
(h(u),e,h(v'))\in V'(S),
\]
using the notation of \SGARefLink{sga3:xxv:lemma:2-8}{2.8}. If the pair on the left is a section of
\(V_1\), then \(h'(h(u),e,h(v'))\) is defined and is a section of
\(\Omega\); write
\[
 h'(h(u),e,h(v'))=(v'',t'',u'').
\]
Set
\[
 \pi\bigl((v,t,u),(v',t',u')\bigr)
 =
 \bigl(v\cdot t v''t^{-1},\ tt''t',\
 t'^{-1}u''t'\cdot u'\bigr).
\]
Assertion (i) follows immediately from \SGARefLink{sga3:xxv:lemma:2-8:item:ii}{2.8(ii)}. To verify the assertion
about \(\pi\) in (iii), observe that
\[
 h'(h(u),e,h(v'))=uv'=v''t''u''
\]
whenever \(u\in U(S)\), \(v'\in U^-(S)\), and \(S\to k\), by
\SGARefLink{sga3:xxv:lemma:2-8:item:iii}{2.8(iii)} and (ii).

The morphism \(\sigma\) is constructed similarly. If \((v,t,u)\) is a
section of \(\Omega\), then \(h(u^{-1})\) is a section of \(U^-\),
\(h(v^{-1})\) a section of \(U\), and \(h(t^{-1})\) a section of \(T\).
Thus
\[
 (h(u^{-1}),h(t^{-1}),h(v^{-1}))
\]
is a section of \(\Omega\).
Define an open subset \(V_2\) of \(\Omega\) by
\[
 (v,t,u)\in V_2(S)
 \quad\Longleftrightarrow\quad
 (h(u^{-1}),h(t^{-1}),h(v^{-1}))\in V'(S),
\]
and define
\[
 \sigma(v,t,u)=h'(h(u^{-1}),h(t^{-1}),h(v^{-1})).
\]
The required properties of \(\sigma\) are verified in the same way.
